% Anexos del Plan. Cada anexo está en capitulos/anexos/; este archivo controla numeración y orientación de página.
\renewcommand{\thetable}{A.\arabic{table}}
\renewcommand{\thefigure}{A.\arabic{figure}}
\renewcommand{\theHtable}{anexo.\arabic{table}}
\renewcommand{\theHfigure}{anexo.\arabic{figure}}
\renewcommand{\theHsection}{anexo.\arabic{section}}
\setcounter{table}{0}
\setcounter{figure}{0}
\setcounter{section}{0}
\begin{singlespace}
\tesisbodyfont
\renewcommand{\thesection}{ANEXO \arabic{section}}
\titleformat{\section}
  {\normalfont\bfseries\fontsize{11}{13.2}\selectfont}
  {\thesection.}{0.75em}{\MakeUppercase}
\titlespacing*{\section}{0pt}{1.0\baselineskip}{0.25\baselineskip}
\begin{landscape}
\chapter*{ANEXOS}
\addcontentsline{toc}{chapter}{ANEXOS}
\thispagestyle{empty}

Los anexos reúnen matrices, diseños y registros que sostienen la trazabilidad del plan. Así respaldan la relación entre problema, objetivos, método, evidencia y evaluación.

\section{MATRIZ DE CONSISTENCIA}
\label{anx:matriz-consistencia}
\thispagestyle{empty}
\scriptsize
La matriz articula los problemas de investigación con los objetivos, productos verificables, insumos, resultados, hipótesis y evidencia de evaluación del artefacto. Los problemas se formulan como enunciados narrativos, de modo que cada fila exprese una carencia verificable que será atendida por el diseño, implementación, evaluación y comunicación del artefacto PINN 2D.

La Tabla~\ref{tab:matriz-consistencia-a} presenta la primera parte de la matriz y verifica la correspondencia entre cada fase DSR, el problema que atiende, el objetivo asociado y el producto que permitirá comprobar su avance.

% ---------------------------------------------------------------------------
% Tabla 1: Desde fase/actividad hasta producto verificable
% ---------------------------------------------------------------------------
\setlength{\tabcolsep}{3pt}
\renewcommand{\arraystretch}{1.18}

\begin{longtable}{%
L{3.8cm}
L{6.2cm}
L{5.3cm}
L{5.0cm}}
\caption{Matriz de consistencia del plan de tesis: problema, objetivo y producto verificable. Fuente: Elaboración propia.}
\label{tab:matriz-consistencia-a}\\
\toprule
\textbf{Fase DSR / actividad} &
\textbf{Problema} &
\textbf{Objetivo} &
\textbf{Producto verificable} \\
\midrule
\endfirsthead

\toprule
\textbf{Fase DSR / actividad} &
\textbf{Problema} &
\textbf{Objetivo} &
\textbf{Producto verificable} \\
\midrule
\endhead

\midrule
\multicolumn{4}{r}{\small Continúa en la siguiente página.}\\
\midrule
\endfoot

\bottomrule
\endlastfoot

\textbf{F0. Definición del objeto de estudio, caso de uso y alcance del artefacto} &
\textbf{PG.} En el arreglo de cables--entorno térmico, la heterogeneidad espacial del suelo nativo, el lecho de asiento y el suelo de relleno puede modificar $T(x,y)$, $\Tmax$ e $\Imax$ respecto de una representación homogénea equivalente \parencites[4--6]{enescu2021}[3--4,93--99]{cigre2025}. Sin embargo, resulta difícil cuantificar esas diferencias de forma reproducible y establecer cuándo son relevantes para las decisiones de diseño y operación. Las representaciones homogéneas equivalentes son rápidas, aunque no describen de forma explícita la conductividad variable $\kx$; los modelos FEM sí pueden representarla, pero exigen construir, mallar y verificar cada escenario. En consecuencia, persiste la necesidad de una formulación computacional verificada y trazable que represente $\kx$, estime el campo térmico y documente la exactitud, la consistencia física y la robustez del análisis. &
\textbf{OG.} Diseñar y evaluar un artefacto computacional basado en una red neuronal informada por física (PINN) que represente el arreglo de cables--entorno térmico, estime $T(x,y)$, $\Tmax$ e $\Imax$, cuantifique el efecto de la heterogeneidad térmica espacial y documente su exactitud, consistencia física, robustez y reproducibilidad. &
Protocolo general del artefacto; dominio de validez; criterios de aceptación; repositorio de escenarios y resultados. \\

\midrule
\textbf{F1. Especificación física, datos y requisitos} &
\textbf{PE1.} La representación física requerida por el artefacto no está formalizada en términos de entradas, requisitos funcionales y criterios de calidad que ordenen el cable multicapa, las interfaces, las fuentes, los contornos y $\kx$ en escenarios heterogéneos. &
\textbf{OE1.} Elaborar una especificación trazable de las entradas físicas, los requisitos funcionales y los criterios de calidad del artefacto para representar el arreglo de cables--entorno térmico, incluyendo dominios multicapa, fuentes térmicas, condiciones de contorno, interfaces y conductividad térmica espacialmente variable $\kx$. &
\textbf{VD1.} Especificación trazable: documento de requisitos; catálogo de materiales y geometrías; ficha de supuestos; matriz requisito--prueba; estado incompleto, completo pendiente o trazable aceptado. \\

\midrule
\textbf{F2. Preparación de datos sintéticos, geometría y referencias} &
\textbf{PE2.} No se dispone de un paquete reproducible de soluciones analíticas o manufacturadas, referencias FEM con convergencia y casos publicados que permita construir y verificar el artefacto dentro del dominio delimitado. &
\textbf{OE2.} Construir y verificar el artefacto a partir de la especificación trazable y de un conjunto de soluciones analíticas o manufacturadas, benchmarks FEM con convergencia y casos publicados, clasificando su estado de aceptación según la exactitud, la consistencia física y la robustez entre ejecuciones. &
\textbf{VD2.} Artefacto verificado: versión del prototipo, base de benchmarks, pruebas, métricas, dominio de validez e informe con estado rechazado, aceptado con restricciones o aceptado. \\

\midrule
\textbf{F3. Experimentación numérica y análisis de sensibilidad} &
\textbf{PE3.} No se ha cuantificado de forma trazable cómo el contraste, el patrón, la extensión y la proximidad de las heterogeneidades modifican $\Tmax$, el campo $T(x,y)$ y la localización del punto caliente. &
\textbf{OE3.} Aplicar el artefacto verificado a un catálogo de escenarios controlados de contraste, patrón, extensión y proximidad de heterogeneidades del suelo nativo y de los suelos de relleno, y consolidar un expediente de evidencia con los cambios en $\Tmax$, el campo de temperatura $T(x,y)$ y la localización del punto caliente. &
\textbf{VD3.} Expediente de evidencia térmica: plan de casos; mapas de $\kx$ y $T(x,y)$; resultados de sensibilidad; punto caliente; límites; y estado inconcluso, sin efecto relevante, con efecto relevante no crítico o con condición crítica. \\

\midrule
\textbf{F4. Evaluación de ampacidad y comparación de escenarios} &
\textbf{PE4.} No se han identificado las condiciones que producen la mayor discrepancia de ampacidad entre una representación 2D heterogénea verificada y una representación homogénea equivalente. &
\textbf{OE4.} Usar el expediente de evidencia térmica y un catálogo de representaciones pareadas para elaborar un expediente comparativo de ampacidad entre la representación 2D heterogénea y las representaciones homogéneas, e identificar las condiciones que generan mayor discrepancia. &
\textbf{VD4.} Expediente comparativo de ampacidad: rutina de búsqueda de $\Imax$; tabla de reducción o aumento; clasificación de escenarios críticos; y estado sin discrepancia relevante, reducción, aumento o inconcluso/fuera del dominio. \\

\end{longtable}

La matriz de consistencia concluye en PE4--OE4 porque la síntesis y la comunicación no expresan una brecha específica adicional. El repositorio depurado, la guía de reproducción, el procedimiento y los principios de diseño se obtendrán en la actividad metodológica descrita en la Sección~\ref{subsec:sintesis-dsr}, a partir de la evidencia producida por las cuatro correspondencias problema--objetivo.

\vspace{0.6cm}

La Tabla~\ref{tab:matriz-consistencia-b} completa la matriz con las variables independientes discretas que ingresan a cada etapa, sus variables dependientes o productos, los estados posibles, las hipótesis, los indicadores y los parámetros controlados. Su función es comprobar que cada VD intermedia se conserva como parte heredada de la VI siguiente, mientras $\kx$, el contraste y los demás atributos físicos permanecen en una columna metodológica distinta.

% ---------------------------------------------------------------------------
% Tabla 2: Entradas discretas, salidas, hipótesis, evidencia y controles
% ---------------------------------------------------------------------------
\setlength{\tabcolsep}{1pt}
\renewcommand{\arraystretch}{1.15}

\begin{longtable}{%
L{2.25cm}
L{3.20cm}
L{3.30cm}
L{3.10cm}
L{4.10cm}
L{3.40cm}
L{3.40cm}}
\caption{Matriz de consistencia del plan de tesis: entradas discretas, salidas, hipótesis, indicadores y parámetros de caso. Fuente: Elaboración propia.}
\label{tab:matriz-consistencia-b}\\
\toprule
\textbf{Fase / objetivo} &
\textbf{VI discreta: composición} &
\textbf{Realizaciones de la VI} &
\textbf{VD discreta: objeto y estados} &
\textbf{Hipótesis} &
\textbf{Indicadores de evaluación} &
\textbf{Parámetros, controles y metadatos} \\
\midrule
\endfirsthead

\toprule
\textbf{Fase / objetivo} &
\textbf{VI discreta: composición} &
\textbf{Realizaciones de la VI} &
\textbf{VD discreta: objeto y estados} &
\textbf{Hipótesis} &
\textbf{Indicadores de evaluación} &
\textbf{Parámetros, controles y metadatos} \\
\midrule
\endhead

\midrule
\multicolumn{7}{r}{\small Continúa en la siguiente página.}\\
\midrule
\endfoot

\bottomrule
\endlastfoot

\textbf{F1 / OE1. Especificación física, datos y requisitos} &
\textbf{VI1=$P_1$. Paquete inicial de requisitos y datos.} No posee entrada heredada. &
Básico; completo; completo documentado. Incluye, según realización, geometría, materiales, fuentes, contornos, interfaces, unidades, supuestos y procedencia. &
\textbf{VD1. Especificación trazable:} requisitos verificables, catálogo de datos, criterios y matriz requisito--prueba. Estados: incompleta; completa pendiente de trazabilidad; trazable aceptada. &
\textbf{HE1.} Un paquete inicial completo y documentado permitirá producir una especificación trazable aceptada, cubrir los requisitos críticos y vincular cada requisito con al menos una prueba física o funcional. &
Porcentaje de requisitos críticos cubiertos; número de pruebas asociadas por requisito; trazabilidad completa/incompleta; cumplimiento de unidades y supuestos. &
Parámetros del caso contenidos en el paquete: geometría, materiales, $\kx$, fuentes y contornos. Controles: norma, sistema de unidades, alcance y exclusiones. \\

\midrule
\textbf{F2 / OE2. Construcción y verificación} &
\textbf{VI2=$\{\mathrm{VD1},R_2\}$. Expediente de construcción y verificación.} Hereda VD1 y agrega referencias. &
VD1 trazable aceptada más $R_2$ analítico/manufacturado; FEM convergente; publicado; o mixto. Sin VD1 aceptada no se habilita OE2. &
\textbf{VD2. Artefacto verificado y expediente de aceptación:} versión, pruebas, métricas y dominio. Estados: rechazado; aceptado con restricciones; aceptado en el dominio. &
\textbf{HE2.} La especificación trazable y referencias válidas permitirán producir un artefacto aceptado dentro del dominio, con error de $\Tmax$ y NRMSE no mayores a 5\,\pct{}, desequilibrio no mayor a 2\,\pct{} y variabilidad menor que el efecto relevante. &
$e_{\Tmax}$; MAE; RMSE/NRMSE; error máximo; residuo PDE; error de contorno; error de interfaz; balance de energía; CV de $\Tmax$. &
Parámetros físicos coincidentes entre artefacto y referencia. Controles/metadatos: dominio, malla y tolerancia FEM, puntos de colocación, arquitectura aceptada, pesos de pérdida y semilla. \\

\midrule
\textbf{F3 / OE3. Evaluación de escenarios} &
\textbf{VI3=$\{\mathrm{VD2},E_3\}$. Expediente de evaluación térmica.} Hereda VD2 y agrega escenarios. &
VD2 aceptada o aceptada con restricciones compatibles más $E_3$ base; estratificado; localizado; mejorado; seco; o multimaterial. &
\textbf{VD3. Expediente de evidencia térmica:} mapas, diferencias y límites. Estados: inconcluso; sin efecto relevante; efecto relevante no crítico; condición crítica. &
\textbf{HE3.} El artefacto verificado aplicado a $E_3$ producirá un expediente concluyente: una región de baja $k$ cercana aumentará $\Tmax$ y una región de alta $k$ usada como relleno la reducirá. &
$\Delta\Tmax$; coordenadas del punto caliente $(x_h,y_h)$; $|\nabla T|_{\max}$; cambio de temperatura media; tamaño de efecto; intervalos por repetición. &
Parámetros físicos del escenario, no VI: mapa $\kx$, $C_k$, $f_h$ y $d_h/D$. Controles: corriente o pérdidas, profundidad, disposición, $T_{\mathrm{amb}}$, dominio, contornos y configuración aceptada. \\

\midrule
\textbf{F4 / OE4. Comparación de ampacidad} &
\textbf{VI4=$\{\mathrm{VD3},C_4\}$. Expediente de comparación de ampacidad.} Hereda VD3 y agrega representaciones pareadas. &
VD3 concluyente más $C_4$: heterogénea 2D frente a homogénea equivalente; conservadora; u optimista. &
\textbf{VD4. Expediente comparativo de ampacidad:} resultados, márgenes y casos críticos. Estados: sin discrepancia relevante; reducción; aumento; inconcluso/fuera del dominio. &
\textbf{HE4.} El expediente térmico combinado con $C_4$ producirá un expediente comparativo en el que la aproximación homogénea tenderá a sobreestimar $\Imax$ al promediar una zona local de baja $k$. &
$\Imax$; $\Delta I\pct$; margen respecto de $T_{\mathrm{lim}}$; número de iteraciones de búsqueda; incertidumbre combinada modelo--referencia. &
Parámetros físicos contenidos en el escenario $E_3$ heredado mediante VD3. Controles: $T_{\mathrm{lim}}$, modelo de pérdidas dependiente de corriente, paso/tolerancia de bisección, geometría y condiciones de borde. \\

\end{longtable}

\end{landscape}

\clearpage
\tesisbodyfont

\begin{landscape}
\section{MATRIZ DE OPERACIONALIZACIÓN DE VARIABLES}
\label{anx:operacionalizacion-variables}
\thispagestyle{empty}

La Tabla~\ref{tab:operacionalizacion-variables} presenta la matriz de operacionalización por etapa.

Este anexo operacionaliza la secuencia entrada--proceso--salida definida en la Sección~\ref{sec:dimensiones-operativas}. Las cuatro variables independientes son paquetes discretos de objetos que ingresan a OE1--OE4; las variables dependientes son los productos discretos de esas etapas. VD1, VD2 y VD3 se conservan como componentes heredados de VI2, VI3 y VI4, respectivamente. Los parámetros físicos, los controles numéricos, los indicadores y los metadatos se registran por separado para impedir que se confundan con las variables de construcción del artefacto.

\footnotesize
\setlength{\tabcolsep}{3pt}
\renewcommand{\arraystretch}{1.18}
\begin{longtable}{L{4.00cm}L{4.20cm}L{8.20cm}L{6.30cm}}
\caption{Matriz de operacionalización de variables. Fuente: Elaboración propia.}
\label{tab:operacionalizacion-variables}\\
\toprule
\textbf{Elemento metodológico} &
\textbf{Variable o elemento} &
\textbf{Operacionalización en el estudio} &
\textbf{Indicador, instrumento o registro} \\
\midrule
\endfirsthead

\toprule
\textbf{Elemento metodológico} &
\textbf{Variable o elemento} &
\textbf{Operacionalización en el estudio} &
\textbf{Indicador, instrumento o registro} \\
\midrule
\endhead

\midrule
\multicolumn{4}{r}{\small Continúa en la siguiente página.}\\
\midrule
\endfoot

\bottomrule
\endlastfoot

Objeto de estudio & Arreglo de cables--entorno térmico & Conjunto físico tangible formado por los cables multicapa y su disposición, el lecho de asiento, el suelo de relleno y el suelo nativo, donde el calor generado en los conductores se transfiere hacia el entorno. & Ficha del caso; geometría y disposición; materiales; carga; condiciones de borde. \\

Unidad de análisis & Sección transversal 2D del objeto de estudio & Caso computacional parametrizado por geometría, materiales, fuentes, contornos, referencia y versión del artefacto. & Identificador de caso; archivo de entrada; metadatos de ejecución; salida reproducible. \\

Variable independiente discreta de OE1 & \textbf{VI1=$P_1$. Paquete inicial de requisitos y datos} & Objeto documental sin entrada heredada. Realizaciones: básico; completo; completo documentado. Define los datos, requisitos y fuentes disponibles para elaborar VD1. & Código de realización; lista de campos; fuentes; versión; ficha de integridad. \\

Variable independiente discreta de OE2 & \textbf{VI2=$\{\mathrm{VD1},R_2\}$. Expediente de construcción y verificación} & Paquete compuesto por VD1 trazable aceptada y un conjunto nuevo $R_2$. Realizaciones de $R_2$: analítica/manufacturada; FEM convergente; publicada; mixta. & Identificador de VD1; código de $R_2$; procedencia; versión; evidencia de convergencia o validez. \\

Variable independiente discreta de OE3 & \textbf{VI3=$\{\mathrm{VD2},E_3\}$. Expediente de evaluación térmica} & Paquete compuesto por VD2 aceptada, o aceptada con restricciones compatibles, y un catálogo nuevo $E_3$. Realizaciones: base; estratificada; localizada; mejorada; seca; multimaterial. & Identificador de VD2; estado heredado; código de $E_3$; archivo de entrada; ficha física. \\

Variable independiente discreta de OE4 & \textbf{VI4=$\{\mathrm{VD3},C_4\}$. Expediente de comparación de ampacidad} & Paquete compuesto por VD3 concluyente y un catálogo nuevo $C_4$. Realizaciones: par heterogéneo 2D--homogéneo equivalente; conservador; optimista. & Identificador de VD3; estado heredado; código de $C_4$; regla de equivalencia; identificador del par. \\

Variable dependiente de OE1 / entrada heredada de OE2 & \textbf{VD1$\equiv\mathrm{VI}_{2,h}$. Especificación trazable} & Producto documental con requisitos, criterios y correspondencia requisito--prueba. Estados: incompleta; completa pendiente de trazabilidad; trazable aceptada. Solo el último estado habilita VI2. & Estado de VD1; cobertura; trazabilidad; consistencia de unidades; supuestos y exclusiones. \\

Variable dependiente de OE2 / entrada heredada de OE3 & \textbf{VD2$\equiv\mathrm{VI}_{3,h}$. Artefacto verificado y expediente de aceptación} & Objeto versionado formado por prototipo, pruebas, métricas y dominio. Estados: rechazado; aceptado con restricciones; aceptado en el dominio. Los estados aceptados compatibles habilitan VI3. & Estado de VD2; error de $\Tmax$; MAE; NRMSE; residuos; balance; robustez; dominio de validez. \\

Variable dependiente de OE3 / entrada heredada de OE4 & \textbf{VD3$\equiv\mathrm{VI}_{4,h}$. Expediente de evidencia térmica} & Objeto documental y de datos por escenario. Estados: inconcluso; sin efecto relevante; efecto relevante no crítico; condición crítica. Los estados concluyentes habilitan VI4. & Estado de VD3; $T(x,y)$; $\Tmax$; $\Delta\Tmax$; gradiente; punto caliente; cobertura del catálogo. \\

Variable dependiente final de OE4 & \textbf{VD4. Expediente comparativo de ampacidad} & Objeto final con resultados, márgenes y casos críticos. Estados: sin discrepancia relevante; reducción de ampacidad; aumento de ampacidad; inconcluso/fuera del dominio. & Estado de VD4; $\Imax$; $\Delta I\pct$; margen; cobertura de pares; clasificación de casos críticos. \\

Parámetros físicos de caso (no son VI ni VD) & Propiedades y configuración del sistema físico & Atributos incorporados en $P_1$, $E_3$ y $C_4$ para describir la física: no representan objetos de transición entre etapas. & Mapa $\kx$; contraste $C_k$; extensión $f_h$; proximidad $d_h/D$; geometría; materiales; pérdidas; $T_{\mathrm{amb}}$; dominio y contornos. \\

Criterios transversales DSR & Exactitud, consistencia física, robustez, trazabilidad y reproducibilidad & Criterios aplicados a VD1--VD4 y a la síntesis metodológica posterior. No constituyen una quinta VD ni un problema específico adicional. & Cumplimiento de umbrales; casos aceptados reproducidos; cobertura documental; límites de validez; bitácora de decisiones. \\

Controles numéricos y metadatos & Configuración aceptada y registro de ejecución & Condiciones que se fijan tras el piloto o se registran para interpretar la evidencia sin atribuirles el efecto físico. & Arquitectura; pesos de pérdida; densidad de puntos; semilla; convergencia; tolerancia FEM; versión de código; tiempo, memoria y hardware. \\

Condiciones no controladas registradas & Incertidumbre documental y variabilidad numérica & Condiciones que no siempre pueden fijarse, pero que se reportan para interpretar la evidencia sin ocultar limitaciones. & Semilla; convergencia; tolerancia FEM; datos faltantes de publicaciones; rango de propiedades térmicas; hardware. \\

Instrumento de registro y evaluación & Artefacto PINN y protocolo de verificación & Conjunto formado por especificación, implementación, casos de referencia, programas de evaluación, bitácora y repositorio. & Código versionado; archivos parametrizados; pruebas; tablas de resultados; figuras térmicas; bitácora DSR. \\

Característica de registro & Registro computacional del campo térmico y métricas escalares & Registro de matrices espaciales, máximos, diferencias porcentuales, residuos, tiempos y metadatos. & Matrices $T(x,y)$; mapas $\kx$; CSV de métricas; reporte de energía; salida gráfica. \\

Tamaño de muestra metodológico & Paquete de referencias, escenarios y pares & Selección intencional de realizaciones de $R_2$, $E_3$ y $C_4$, acompañadas por los productos heredados aceptados y por combinaciones de parámetros físicos suficientes para verificar el artefacto y cubrir escenarios relevantes. & Casos analíticos o manufacturados; referencias FEM; casos publicados; malla de cobertura del \ref{anx:diseno-casos}. \\

\end{longtable}

\end{landscape}

\tesisbodyfont

\clearpage
\begin{landscape}
\section{DISEÑO DISCRETO DE CASOS Y PLAN DE COBERTURA}
\label{anx:diseno-casos}
\thispagestyle{empty}

La Tabla~\ref{tab:diseno-casos} presenta el diseño discreto de casos y los parámetros físicos de cobertura.

Este anexo define cómo se formarán VI3 y VI4 y cómo se cubrirá el espacio físico de los casos. VI3 combina el artefacto verificado VD2 con el catálogo de escenarios $E_3$; VI4 combina el expediente térmico VD3 con el catálogo de representaciones $C_4$. El contraste, la extensión y la proximidad son parámetros que caracterizan un escenario ya seleccionado; por ello, sus niveles forman una malla de sensibilidad física y no un diseño factorial de variables independientes.

\scriptsize
\setlength{\tabcolsep}{3pt}
\renewcommand{\arraystretch}{1.18}
\begin{longtable}{L{3.10cm}L{1.45cm}L{3.00cm}L{5.15cm}L{5.25cm}}
\caption{Diseño discreto de casos y parámetros físicos de cobertura. Fuente: Elaboración propia.}
\label{tab:diseno-casos}\\
\toprule
\textbf{Elemento} &
\textbf{Código} &
\textbf{Papel metodológico} &
\textbf{Categorías o niveles iniciales} &
\textbf{Uso en el plan de cobertura} \\
\midrule
\endfirsthead

\toprule
\textbf{Elemento} &
\textbf{Código} &
\textbf{Papel metodológico} &
\textbf{Categorías o niveles iniciales} &
\textbf{Uso en el plan de cobertura} \\
\midrule
\endhead

\midrule
\multicolumn{5}{r}{\small Continúa en la siguiente página.}\\
\midrule
\endfoot

\bottomrule
\endlastfoot

Artefacto verificado heredado & VD2 / $\mathrm{VI}_{3,h}$ & Producto de OE2 y componente obligatorio de VI3 & Aceptado en el dominio; aceptado con restricciones compatibles. El estado rechazado no pasa a OE3. & Fija la versión del artefacto y su dominio de validez antes de recorrer escenarios. \\

Catálogo nuevo de escenarios & $E_3$ & Componente nuevo de VI3 & Base homogéneo; estratificado; inclusión localizada; material de mejora; zona seca; combinación multimaterial. & Selecciona el objeto físico-computacional; cada realización contiene mapa $\kx$ y ficha de parámetros. \\

Expediente de evaluación térmica & VI3 & Variable independiente discreta de OE3: $\{\mathrm{VD2},E_3\}$ & Combinaciones válidas de una VD2 aceptada con una realización de $E_3$. & Constituye la entrada completa de OE3 y produce VD3. \\

Expediente térmico heredado & VD3 / $\mathrm{VI}_{4,h}$ & Producto de OE3 y componente obligatorio de VI4 & Sin efecto relevante; efecto relevante no crítico; condición crítica. El estado inconcluso no pasa a OE4. & Conserva mapas, métricas, límites y clasificación del escenario que será comparado. \\

Catálogo nuevo de representaciones & $C_4$ & Componente nuevo de VI4 & Par heterogéneo 2D--homogéneo equivalente; conservador; optimista. & Define la regla de equivalencia y el par que se usará para calcular la discrepancia de ampacidad. \\

Expediente de comparación & VI4 & Variable independiente discreta de OE4: $\{\mathrm{VD3},C_4\}$ & Combinaciones válidas de una VD3 concluyente con una realización de $C_4$. & Constituye la entrada completa de OE4 y produce VD4. \\

Contraste térmico & $C_k$ & Parámetro físico de caso; no VI ni VD & Bajo; medio; alto, con magnitudes fijadas en el piloto para cada realización de $E_3$. & Cubre diferencias de conductividad entre la región heterogénea y el entorno base y permite interpretar los indicadores contenidos en VD3 y VD4. \\

Extensión de la región heterogénea & $f_h$ & Parámetro físico de caso; no VI ni VD & Pequeña; media; grande, expresada como fracción del entorno local o tamaño normalizado. & Cubre el tamaño de la barrera o mejora térmica dentro de cada realización de $E_3$. \\

Proximidad al conductor & $d_h/D$ & Parámetro físico de caso; no VI ni VD & Cercana; media; lejana, normalizada por el diámetro exterior del cable o una distancia característica $D$. & Cubre la ubicación de la región respecto de la zona de mayor flujo térmico. \\

Condiciones pareadas & --- & Controles; no VI & Geometría, pérdidas, contornos, dominio, $T_{\mathrm{amb}}$ y configuración numérica aceptada. & Se mantienen constantes dentro de cada comparación para interpretar el efecto del atributo físico o de la representación. \\

Semilla o repetición numérica & $S$ & Metadato de robustez; no VI & Repeticiones seleccionadas en el piloto y en casos extremos. & Cuantifica variabilidad numérica de las salidas sin convertir el entrenamiento en una variable del estudio. \\

\end{longtable}

\tesisbodyfont

La malla máxima de cobertura combinará tres niveles de $C_k$, $f_h$ y $d_h/D$ dentro de cada una de las cuatro realizaciones heterogéneas simples de $E_3$. Esto produce hasta $3^3=27$ configuraciones paramétricas por realización.

Si se recorren completas esas cuatro realizaciones, el máximo de planificación será de 108 escenarios heterogéneos. Cada VD3 concluyente se combinará con la realización de $C_4$ que corresponda para formar VI4 y construir su comparación homogénea. El número 108 describe cobertura paramétrica potencial; no es el número de estados de VI3 o VI4 ni implica inferencia estadística factorial.

Este número se entiende como marco de planificación. Por tanto, si la carga computacional del piloto supera los recursos disponibles, se conservarán completos los patrones más críticos y se justificará cualquier fracción usada.

La lectura de resultados se hará en tres capas. Primero, se compararán los indicadores contenidos en VD3 y VD4 entre realizaciones discretas de $E_3$ y $C_4$, y se asignará a cada expediente uno de sus estados previstos.

Segundo, se describirá la sensibilidad física dentro de cada realización de $E_3$, especialmente para las combinaciones $C_k\times d_h/D$ y $C_k\times f_h$. Estas expresiones identifican cruces de parámetros del caso, no interacciones entre variables independientes metodológicas.

Tercero, se contrastará cada escenario heterogéneo con su homogéneo equivalente. Así se identificará cuándo la simplificación introduce una diferencia térmica u operativa relevante.

Para mantener consistencia con la matriz de operacionalización, las comparaciones pareadas mantendrán fijos geometría, pérdidas, contornos, dominio y configuración numérica. Por tanto, el efecto observado se atribuirá al atributo físico o al tipo de representación que diferencia el par.

Las condiciones no controladas se registrarán como metadatos. Entre ellas estarán la semilla, la tolerancia FEM publicada y la incertidumbre de las propiedades térmicas.


\end{landscape}
\clearpage
\tesisbodyfont

% El anexo 4 alterna vertical y horizontal internamente.
\section{ESTRATEGIA DE REVISIÓN DOCUMENTAL}
\thispagestyle{empty}

La base documental se construirá con fuentes académicas, normativas y técnicas reconocidas. Además, se usarán fuentes recientes o vigentes, recuperadas como texto completo y organizadas en Zotero.

Para los antecedentes se priorizarán artículos revisados por pares de 2021--2026. Sin embargo, no se excluirán normas vigentes ni fuentes clásicas indispensables.

Las búsquedas combinarán términos de cable subterráneo, ampacidad, suelo nativo, suelo de relleno, heterogeneidad, FEM, PINN, conducción de calor y DSR. Por tanto, cada afirmación cuantitativa se vinculará con página, tabla o figura de la fuente primaria.

La estrategia se ejecutó el 5 de junio de 2026. Además de auditar los registros del gestor bibliográfico y los PDF locales, se verificaron DOI, títulos y disponibilidad de texto completo mediante fuentes editoriales, Crossref y OpenAlex. El string general de partida fue:

\begin{quote}\scriptsize\ttfamily\RaggedRight
("underground power cable*" OR "buried power cable*" OR "underground cable system*")\\
AND (thermal OR ampacity OR "thermal rating" OR temperature)\\
AND (review OR assessment OR systematic OR survey OR mapping OR overview)
\end{quote}

Para mantener la búsqueda reproducible sin extender este anexo, la Tabla~\ref{tab:strings-documentales} se incluye como evidencia operativa de la estrategia y resume los bloques booleanos usados por tema. En cada bloque se aplicaron, según el tipo de fuente requerida, los filtros estado del arte: \texttt{review OR systematic OR bibliometric OR survey OR mapping OR overview OR assessment}; aporte: \texttt{method OR framework OR model OR formulation OR algorithm OR "inverse problem" OR "calculation method"}; y aplicación: \texttt{application OR prediction OR reconstruction OR monitoring OR simulation OR experiment OR optimization OR "case study"}.

La Tabla~\ref{tab:criterios-documentales} se incluye para explicitar los criterios de selección documental y reducir la arbitrariedad al aceptar, excluir o usar una fuente solo como apoyo.

\clearpage
\begin{landscape}
\thispagestyle{empty}
\begin{center}
\captionsetup{type=table,hypcap=false}
\phantomsection
\caption{Strings de búsqueda documental usados por tema. Fuente: Elaboración propia.}
\label{tab:strings-documentales}
\vspace{0.4em}
\footnotesize
\setlength{\tabcolsep}{5pt}
\renewcommand{\arraystretch}{1.22}
\begin{tabularx}{\linewidth}{>{\bfseries}L{4.1cm}>{\ttfamily\scriptsize\RaggedRight\arraybackslash}X}
\toprule
\normalfont\bfseries Tema & \normalfont\bfseries Bloque booleano usado \\
\midrule
PINNs & ("physics-informed neural network*" OR PINN* OR XPINN* OR "scientific machine learning" OR "physics-enhanced" OR "FEM-BPNN") AND ("partial differential equation*" OR PDE OR "heat conduction" OR "heat transfer" OR "thermal conductivity" OR "underground power cable*" OR "temperature field") \\
\addlinespace[0.35em]
Transferencia de calor & ("heat transfer" OR "heat conduction" OR "thermal assessment" OR "thermal model" OR "temperature field") AND ("power cable*" OR "underground cable*" OR "buried cable*" OR soil OR backfill OR "multi-layered soil") \\
\addlinespace[0.35em]
Suelo nativo, suelo de relleno y humedad & ("soil thermal conductivity" OR "soil thermal resistivity" OR "thermal backfill" OR "soil moisture" OR "unsaturated soils") AND ("underground cable*" OR "buried infrastructure" OR "power cable*" OR measurement OR model* OR moisture OR backfill) \\
\addlinespace[0.35em]
FEM y métodos numéricos & ("finite element method" OR FEM OR "finite element" OR "finite difference" OR "numerical simulation" OR "spectral finite element" OR "adaptive finite element") AND ("heat transfer" OR "heat conduction" OR "heterogeneous media" OR "multi-layered soil" OR "power cable rating") \\
\addlinespace[0.35em]
Cables y ampacidad & ("power cable*" OR "underground power cable*" OR "buried power cable*" OR "insulated cable*") AND (ampacity OR "current rating" OR "thermal rating" OR "temperature rise" OR "load capability") \\
\addlinespace[0.35em]
DTR y cargas cíclicas & ("dynamic thermal rating" OR DTR OR "cyclic loading" OR "emergency rating" OR "dynamic thermal analysis" OR "transient temperature") AND ("underground cable*" OR "XLPE cable*" OR "power cable*" OR "transmission line*") \\
\addlinespace[0.35em]
Normas y estándares & (IEC OR IEEE OR CIGRE OR ASTM OR standard*) AND (ampacity OR "current rating" OR "thermal rating" OR "cable ampacity" OR "calculation method" OR "measurement method") AND ("power cable*" OR "underground cable*" OR soil OR backfill) \\
\bottomrule
\end{tabularx}

\vspace{1.0em}

\phantomsection
\caption{Criterios de inclusión y exclusión documental. Fuente: Elaboración propia.}
\label{tab:criterios-documentales}
\footnotesize
\setlength{\tabcolsep}{5pt}
\renewcommand{\arraystretch}{1.18}
\begin{tabularx}{\linewidth}{L{3.2cm}Y Y}
\toprule
\textbf{Criterio} & \textbf{Incluir} & \textbf{Excluir o usar solo como apoyo} \\
\midrule
Pertinencia & Transferencia de calor, ampacidad, heterogeneidad, PINN térmica o DSR. & Inteligencia artificial (IA) o cables sin relación con el problema de investigación. \\
Calidad & Artículo arbitrado, norma, guía técnica o libro reconocido. & Sitios sin autoría o metodología verificable. \\
Actualidad & 2021--2026 para estado del arte. & Fuentes antiguas salvo normas, fundamentos o benchmarks. \\
Trazabilidad & Texto completo y localizador de página. & Resúmenes sin acceso al método o resultado citado. \\
Duplicación & Versión final y registro canónico. & Copias duplicadas de la misma publicación. \\
\bottomrule
\end{tabularx}
\end{center}
\end{landscape}

\clearpage

\end{singlespace}

% El anexo 5 se compone fuera del bloque a interlineado sencillo.
\clearpage
\thispagestyle{plain}
\tesisbodyfont
\renewcommand{\thesection}{ANEXO \arabic{section}}
\titleformat{\section}
  {\normalfont\bfseries\fontsize{11}{13.2}\selectfont}
  {\thesection.}{0.75em}{\MakeUppercase}
\titlespacing*{\section}{0pt}{1.0\baselineskip}{0.25\baselineskip}
% Requiere en el preámbulo:
% \usepackage{hyperref}
% \usepackage{xurl} % opcional, ayuda a cortar URLs largas

\section{REPOSITORIO DIGITAL DEL INFORME}

En la presente hoja se consigna la dirección URL del repositorio digital en Google Drive donde se encuentra la documentación complementaria del informe. Además, se resume el tipo de materiales que respaldan la trazabilidad documental del plan.

\vspace{0.5cm}

\noindent\textbf{Dirección URL del repositorio:}

\vspace{0.3cm}

\noindent
\url{https://drive.google.com/drive/folders/1t_vPJ7YbWgOVJ2CKJ4A22lGIAcpq6uH5?usp=drive_link}

\vspace{0.8cm}

\noindent\textbf{Contenido del repositorio:}

\begin{enumerate}
    \item Papers en formato PDF y sus respectivos resúmenes, con los archivos renombrados según el título del paper.
    
    \item Ejemplos de datos compilados o parametrizados para los casos de estudio.
    
    \item Normas, estándares, guías y manuales relacionados con el tema de investigación.
    
    \item Archivos originales de los documentos entregados en PDF y \LaTeX.
    
    \item Archivo Excel que contiene la matriz de consistencia.
\end{enumerate}

\vfill

\noindent
Este repositorio constituye el respaldo digital de los documentos, fuentes, datos y anexos utilizados para el desarrollo del informe. Por tanto, apoya la trazabilidad de las fuentes y de los materiales complementarios del plan.





